\chapter*{Abstract}
\addcontentsline{toc}{chapter}{Abstract}

\begin{tabular}{@{} >{\bfseries}p{0.2\textwidth} p{0.77\textwidth} @{}}
\texorpdfstring{\MakeUppercase{Keywords}}{Keywords} &
telecom tower sites; hybrid energy systems; solar PV; battery storage;
diesel logistics; reinforcement learning; constrained Markov decision process;
inventory management; action masking; Expected Energy Not Served (EENS);
model predictive control
\end{tabular}
\par
\begin{doublespacing}

Reliable and low-carbon energy management for hybrid telecom tower sites
remains a critical challenge in regions with weak or intermittent electricity
supply. \revised{Conventional rule-based and optimisation-based
controllers offer practical energy-management
approaches, but their performance depends on
appropriate calibration, system assumptions
and the representation of fuel replenishment
uncertainty.}

This thesis studies the operational control of hybrid
solar--battery--diesel telecom tower systems under uncertainty in renewable
generation, grid availability, and diesel delivery. The problem is formulated
as a Constrained Markov Decision Process (CMDP) that explicitly models
battery operating constraints and diesel inventory dynamics, including diesel replenishment ordering and stochastic delivery lead times.

A data-grounded simulation environment is developed using the ITU~2024
Smart Energy Scheduling dataset, comprising operational traces from ten
telecom base station sites at hourly resolution. \revised{The proposed RLInv controller uses hard
action masking to exclude infeasible actions before
selection. Unmasked reinforcement-learning
configurations are evaluated separately as
comparators.}

\begin{revisedblock}
The original evaluation compares RLInv with eight
alternative controller configurations across ten
sites and four diesel-delivery scenarios. The core
learning comparisons use ten independent training
seeds. Supplementary experiments examine calibrated
classical controllers, non-learning reference rules,
alternative reinforcement-learning algorithms,
computational cost, and cross-site generalisation.

Within the original primary baseline comparison,
RLInv achieves the lowest mean Expected Energy
Not Served (EENS) against B0, B1, MPC,
Oracle-MPC and TrackB across all four logistics
scenarios. The separate ablation controllers
are not included in this ranking. However,
the subsequent site-level RLInv--Oracle-MPC analysis
does not establish statistically significant
differences. In the RLInv--A6 ablation, the
difference remains significant under Extreme
logistics after correction for multiple comparisons,
but not under the other scenarios. Separate
ablations find no statistically significant
mean-EENS benefit from action masking or the
combined inventory-observation channels.

The expanded comparisons qualify the initial
results. A calibrated classical policy approaches
or improves on RLInv under stressed logistics,
while a simple non-learning controller combining
greedy diesel dispatch with feasible maximum
replenishment matches or improves on RLInv's mean
EENS across all four scenarios. Under the fixed
dispatch priority, unserved energy and diesel
generation are linked by an energy-conservation
relationship. The implemented reward also omits
explicit replenishment and inventory-holding costs.

Cross-site experiments additionally evaluate
shared-policy transfer to five held-out sites,
with comparable reliability in many matched
comparisons and fewer nominal training interactions.
Overall, the findings establish the relevance
of diesel inventory and replenishment decisions,
while showing that the incremental advantage of
learned control depends on the benchmark,
reward formulation, and reference policies.
\end{revisedblock}

\end{doublespacing}